% Options for packages loaded elsewhere
% Options for packages loaded elsewhere
\PassOptionsToPackage{unicode}{hyperref}
\PassOptionsToPackage{hyphens}{url}
%
\documentclass[
  english,
  russian,
  12pt,
  a4paper,
  DIV=11,
  numbers=noendperiod]{scrreprt}
\usepackage{xcolor}
\usepackage{amsmath,amssymb}
\setcounter{secnumdepth}{5}
\usepackage{iftex}
\ifPDFTeX
  \usepackage[T1]{fontenc}
  \usepackage[utf8]{inputenc}
  \usepackage{textcomp} % provide euro and other symbols
\else % if luatex or xetex
  \usepackage{unicode-math} % this also loads fontspec
  \defaultfontfeatures{Scale=MatchLowercase}
  \defaultfontfeatures[\rmfamily]{Ligatures=TeX,Scale=1}
\fi
\usepackage{lmodern}
\ifPDFTeX\else
  % xetex/luatex font selection
\fi
% Use upquote if available, for straight quotes in verbatim environments
\IfFileExists{upquote.sty}{\usepackage{upquote}}{}
\IfFileExists{microtype.sty}{% use microtype if available
  \usepackage[]{microtype}
  \UseMicrotypeSet[protrusion]{basicmath} % disable protrusion for tt fonts
}{}
\usepackage{setspace}
% Make \paragraph and \subparagraph free-standing
\makeatletter
\ifx\paragraph\undefined\else
  \let\oldparagraph\paragraph
  \renewcommand{\paragraph}{
    \@ifstar
      \xxxParagraphStar
      \xxxParagraphNoStar
  }
  \newcommand{\xxxParagraphStar}[1]{\oldparagraph*{#1}\mbox{}}
  \newcommand{\xxxParagraphNoStar}[1]{\oldparagraph{#1}\mbox{}}
\fi
\ifx\subparagraph\undefined\else
  \let\oldsubparagraph\subparagraph
  \renewcommand{\subparagraph}{
    \@ifstar
      \xxxSubParagraphStar
      \xxxSubParagraphNoStar
  }
  \newcommand{\xxxSubParagraphStar}[1]{\oldsubparagraph*{#1}\mbox{}}
  \newcommand{\xxxSubParagraphNoStar}[1]{\oldsubparagraph{#1}\mbox{}}
\fi
\makeatother


\usepackage{longtable,booktabs,array}
\usepackage{calc} % for calculating minipage widths
% Correct order of tables after \paragraph or \subparagraph
\usepackage{etoolbox}
\makeatletter
\patchcmd\longtable{\par}{\if@noskipsec\mbox{}\fi\par}{}{}
\makeatother
% Allow footnotes in longtable head/foot
\IfFileExists{footnotehyper.sty}{\usepackage{footnotehyper}}{\usepackage{footnote}}
\makesavenoteenv{longtable}
\usepackage{graphicx}
\makeatletter
\newsavebox\pandoc@box
\newcommand*\pandocbounded[1]{% scales image to fit in text height/width
  \sbox\pandoc@box{#1}%
  \Gscale@div\@tempa{\textheight}{\dimexpr\ht\pandoc@box+\dp\pandoc@box\relax}%
  \Gscale@div\@tempb{\linewidth}{\wd\pandoc@box}%
  \ifdim\@tempb\p@<\@tempa\p@\let\@tempa\@tempb\fi% select the smaller of both
  \ifdim\@tempa\p@<\p@\scalebox{\@tempa}{\usebox\pandoc@box}%
  \else\usebox{\pandoc@box}%
  \fi%
}
% Set default figure placement to htbp
\def\fps@figure{htbp}
\makeatother



\ifLuaTeX
\usepackage[bidi=basic,provide=*]{babel}
\else
\usepackage[bidi=default,provide=*]{babel}
\fi
% get rid of language-specific shorthands (see #6817):
\let\LanguageShortHands\languageshorthands
\def\languageshorthands#1{}


\setlength{\emergencystretch}{3em} % prevent overfull lines

\providecommand{\tightlist}{%
  \setlength{\itemsep}{0pt}\setlength{\parskip}{0pt}}



 
\usepackage[style=gost-numeric,backend=biber,langhook=extras,autolang=other*]{biblatex}
\addbibresource{bib/cite.bib}

\usepackage[]{csquotes}

\usepackage{indentfirst}
\usepackage{float}
\floatplacement{figure}{H}
\IfFileExists{plex-otf.sty}{
  %% Full TeXlive
  \usepackage[
    % math,
    RM={Scale=0.94},SS={Scale=0.94},SScon={Scale=0.94},TT={Scale=MatchLowercase,FakeStretch=0.9},
    DefaultFeatures={Ligatures=Common}
  ]{plex-otf}
  % \usepackage[Scale=MatchUppercase]{juliamono}
}{
  %% TinyTeX
  \usepackage{libertine}
}
\KOMAoption{captions}{tableheading}
\makeatletter
\@ifpackageloaded{caption}{}{\usepackage{caption}}
\AtBeginDocument{%
\ifdefined\contentsname
  \renewcommand*\contentsname{Содержание}
\else
  \newcommand\contentsname{Содержание}
\fi
\ifdefined\listfigurename
  \renewcommand*\listfigurename{Список иллюстраций}
\else
  \newcommand\listfigurename{Список иллюстраций}
\fi
\ifdefined\listtablename
  \renewcommand*\listtablename{Список таблиц}
\else
  \newcommand\listtablename{Список таблиц}
\fi
\ifdefined\figurename
  \renewcommand*\figurename{Рисунок}
\else
  \newcommand\figurename{Рисунок}
\fi
\ifdefined\tablename
  \renewcommand*\tablename{Таблица}
\else
  \newcommand\tablename{Таблица}
\fi
}
\@ifpackageloaded{float}{}{\usepackage{float}}
\floatstyle{ruled}
\@ifundefined{c@chapter}{\newfloat{codelisting}{h}{lop}}{\newfloat{codelisting}{h}{lop}[chapter]}
\floatname{codelisting}{Список}
\newcommand*\listoflistings{\listof{codelisting}{Листинги}}
\makeatother
\makeatletter
\makeatother
\makeatletter
\@ifpackageloaded{caption}{}{\usepackage{caption}}
\@ifpackageloaded{subcaption}{}{\usepackage{subcaption}}
\makeatother
\usepackage{bookmark}
\IfFileExists{xurl.sty}{\usepackage{xurl}}{} % add URL line breaks if available
\urlstyle{same}
\hypersetup{
  pdftitle={Индивидуальный проект},
  pdfauthor={Курушин Георгий Романович},
  pdflang={ru-RU},
  hidelinks,
  pdfcreator={LaTeX via pandoc}}


\title{Индивидуальный проект}
\usepackage{etoolbox}
\makeatletter
\providecommand{\subtitle}[1]{% add subtitle to \maketitle
  \apptocmd{\@title}{\par {\large #1 \par}}{}{}
}
\makeatother
\subtitle{4 Этап. Использование nikto}
\author{Курушин Георгий Романович}
\date{}
\begin{document}
\maketitle

\renewcommand*\contentsname{Содержание}
{
\setcounter{tocdepth}{1}
\tableofcontents
}
\listoffigures
\listoftables

\setstretch{1.5}
\chapter{Цель
работы}\label{ux446ux435ux43bux44c-ux440ux430ux431ux43eux442ux44b}

Научиться тестированию веб-приложений с помощью сканера nikto

\chapter{Задание}\label{ux437ux430ux434ux430ux43dux438ux435}

Просканировать веб-приложение, используя nikto

\chapter{Теоретическое
введение}\label{ux442ux435ux43eux440ux435ux442ux438ux447ux435ux441ux43aux43eux435-ux432ux432ux435ux434ux435ux43dux438ux435}

\emph{nikto} --- базовый сканер безопасности веб-сервера. Он сканирует и
обнаруживает уязвимости в веб-приложениях, обычно вызванные неправильной
конфигурацией на самом сервере, файлами, установленными по умолчанию, и
небезопасными файлами, а также устаревшими серверными приложениями.
Поскольку nikto построен исключительно на LibWhisker2, он сразу после
установки поддерживает кросс-платформенное развертывание, SSL
(криптографический протокол, который подразумевает более безопасную
связь), методы аутентификации хоста (NTLM/Basic), прокси и несколько
методов уклонения от идентификаторов. Он также поддерживает перечисление
поддоменов, проверку безопасности приложений (XSS, SQL-инъекции и т. д.)
и способен с помощью атаки паролей на основе словаря угадывать учетные
данные авторизации.

Для запуска сканера nikto введите в командную строку терминала команду:
\texttt{\#\ nikto}

По умолчанию, как ранее было показано в других приложениях, при обычном
запуске команды отображаются различные доступные параметры. Для
сканирования цели введите
\texttt{nikto\ -h\ \textless{}цель\textgreater{}\ -p\ \textless{}порт\textgreater{}},
где --- домен или IP-адрес целевого сайта, а --- порт, на котором
запущен сервис

Сканер nikto позволяет идентифицировать уязвимости веб-приложений, такие
как раскрытие информации, инъекция (XSS/Script/HTML), удаленный поиск
файлов (на уровне сервера), выполнение команд и идентификация
программного обеспечения. В дополнение к показанному ранее основному
сканированию nikto позволяет испытателю на проникновение настроить
сканирование конкретной цели. Рассмотрим параметры, которые следует
использовать при сканировании. - Указав переключатель командной строки
-T с отдельными номерами тестов, можно настроить тестирование конкретных
типов. - Используя при тестировании параметр -t, вы можете установить
значение тайм-аута для каждого ответа. - Параметр -D V управляет выводом
на экран. - Параметры -o и -F отвечают за выбор формата отчета
сканирования.

Существуют и другие параметры, такие как -mutate (угадывать поддомены,
файлы, каталоги и имена пользователей), -evasion (обходить фильтр
идентификаторов) и -Single (для одиночного тестового режима), которые
можно использовать для углубленной оценки цели \autocite{parasram}.

\chapter{Выполнение
работы}\label{ux432ux44bux43fux43eux43bux43dux435ux43dux438ux435-ux440ux430ux431ux43eux442ux44b}

Запускаю mysql и apache2 (рис.~\ref{fig-001}).

\begin{figure}

\centering{

\includegraphics[width=0.7\linewidth,height=\textheight,keepaspectratio]{image/1.png}

}

\caption{\label{fig-001}Запуск mysql и apache2}

\end{figure}%

Ввожу в адресной строке браузера адрес DVWA и перехожу в режим выбора
уровня безопасности, сначала поставлю минимальный - Low
(рис.~\ref{fig-002}).

\begin{figure}

\centering{

\includegraphics[width=0.7\linewidth,height=\textheight,keepaspectratio]{image/2.png}

}

\caption{\label{fig-002}DVWA - выбор уровня безопасности}

\end{figure}%

Запускаю nikto и проверяю веб-приложение, ввелдя сначала адрес хоста и
порта (рис.~\ref{fig-003}).

\begin{figure}

\centering{

\includegraphics[width=0.7\linewidth,height=\textheight,keepaspectratio]{image/3.png}

}

\caption{\label{fig-003}Сканирование через ввод адреса хоста и порта}

\end{figure}%

Теперь пробую просканировать, введя полный URL-адрес
(рис.~\ref{fig-004}). Результаты незначительно отличаются.

\begin{figure}

\centering{

\includegraphics[width=0.7\linewidth,height=\textheight,keepaspectratio]{image/4.png}

}

\caption{\label{fig-004}Сканирование через ввод URL-адреса}

\end{figure}%

Попробую проделать теперь то же самое, изменив уровень безопасности на
impossible (рис.~\ref{fig-005}).

\begin{figure}

\centering{

\includegraphics[width=0.7\linewidth,height=\textheight,keepaspectratio]{image/5.png}

}

\caption{\label{fig-005}Выбор самого высокого уровня безопасности}

\end{figure}%

Повторив предыдущую команду, я не заметил изменений в выводе информации
(рис.~\ref{fig-006}).

\begin{figure}

\centering{

\includegraphics[width=0.7\linewidth,height=\textheight,keepaspectratio]{image/6.png}

}

\caption{\label{fig-006}Сканирование после изменения уровня
безопасности}

\end{figure}%

\section{Анализ результатов
сканирования}\label{ux430ux43dux430ux43bux438ux437-ux440ux435ux437ux443ux43bux44cux442ux430ux442ux43eux432-ux441ux43aux430ux43dux438ux440ux43eux432ux430ux43dux438ux44f}

Кроме адреса хоста и порта веб-приложения, nikto выводит инофрмацию о
различных уязвимостях приложения:

Сервер: Apache/2.4.58 (Debian) + /DVWA/: Заголовок X-Frame-Options,
защищающий от перехвата кликов, отсутствует. Смотрите:
https://developer.mozilla.org/en-US/docs/Web/HTTP/Headers/X-Frame-Options

\begin{itemize}
\item
  /DVWA/: Заголовок X-Content-Type-Options не задан. Это может позволить
  пользовательскому агенту отображать содержимое сайта способом,
  отличным от MIME-типа. Смотрите:
  https://www.netsparker.com/web-vulnerability-scanner/vulnerabilities/missing-content-type-header/
\item
  Корневая страница /DVWA перенаправляет на: login.php
\item
  Каталоги CGI не найдены (используйте \enquote*{-C all}, чтобы
  принудительно проверить все возможные каталоги)
\item
  ОПЦИИ: Разрешенные HTTP-методы: GET, POST, OPTIONS, HEAD .
\item
  /DVWA///etc/hosts: Установка сервера позволяет считывать любой
  системный файл, добавляя дополнительный \enquote{/} к URL-адресу.
\item
  /DVWA/config/: Найдена индексация каталога.
\item
  /DVWA/config/: Информация о конфигурации может быть доступна удаленно.
\item
  /DVWA/tests/: Найдена индексация каталога.
\item
  /DVWA/tests/: Это может быть интересно.
\item
  /DVWA/database/: Найдена индексация каталога.
\item
  /DVWA/база данных/: Найден каталог базы данных.
\item
  /DVWA/документы/: Найдена индексация каталога.
\item
  /DVWA/login.php: Найдена страница входа администратора/раздел.
\item
  /DVWA/.git/index: Индексный файл Git может содержать информацию о
  списке каталогов.
\item
  /DVWA/.git/HEAD: Найден файл Git HEAD. Может содержаться полная
  информация о репозитории.
\item
  /DVWA/.git/config: Найден конфигурационный файл Git. Может содержаться
  информация о деталях репозитория.
\item
  /DVWA/.gitignore: найден файл .gitignore. Можно разобраться в
  структуре каталогов.
\item
  /DVWA/wp-content/themes/twentyeleven/images/headers/server.php?filesrc=/etc/hosts:
  Обнаружен файловый менеджер с бэкдором на PHP.
\item
  /DVWA/wordpress/wp-content/themes/twentyeleven/images/headers/server.php?filesrc=/etc/hosts:
  Обнаружен файловый менеджер с бэкдором на PHP.
\item
  /DVWA/wp-includes/Requests/Utility/content-post.php?filesrc=/etc/hosts:
  Найден файловый менеджер с бэкдором на PHP.
\item
  /DVWA/wordpress/wp-includes/Requests/Utility/content-post.php?filesrc=/etc/hosts:
  Найден файловый менеджер с бэкдором на PHP.
\item
  /DVWA/wp-включает в
  себя/js/tinymce/themes/modern/Meuhy.php?filesrc=/etc/hosts: Найден
  файловый менеджер бэкдора PHP.
\item
  /DVWA/wordpress/wp-включает в
  себя/js/tinymce/themes/modern/Meuhy.php?filesrc=/etc/hosts: Найден
  файловый менеджер бэкдора на PHP.
\item
  /DVWA/assets/mobirise/css/meta.php?filesrc=: Найден файловый менеджер
  бэкдора на PHP.
\item
  /DVWA/login.cgi?cli=aa\%20aa\%27cat\%20/etc/hosts: Удаленное
  выполнение какой-либо команды маршрутизатором D-Link.
\item
  /DVWA/shell?cat+/etc/hosts: Обнаружен черный ход.
\item
  /DVWA/.dockerignore: найден файл .dockerignore. Возможно, удастся
  разобраться в структуре каталогов и узнать больше о сайте.
\end{itemize}

Бэкдор, тайный вход (от англ. back door --- «чёрный ход», «лазейка»,
буквально «задняя дверь») --- дефект алгоритма, который намеренно
встраивается в него разработчиком и позволяет получить
несанкционированный доступ к данным или удалённому управлению
операционной системой и компьютером в целом.

Также в результатах nikto отображает код OSVDB 561 и дает ссылку на
CVE-2003-1418. OSVDB --- это аббревиатура базы данных уязвимостей с
открытым исходным кодом.

CVE-2003-1418 --- это уязвимость в Apache HTTP Server 1.3.22--1.3.27 на
OpenBSD, которая позволяет удалённым злоумышленникам получать
конфиденциальную информацию через:

\begin{itemize}
\item
  Заголовок ETag, который раскрывает номер вode.
\item
  Многочастную границу MIME, которая раскрывает идентификаторы дочерних
  процессов (PID).
\end{itemize}

В настоящее время эта проблема имеет среднюю степень тяжести.

\chapter{Выводы}\label{ux432ux44bux432ux43eux434ux44b}

В ходе работы я научился использовать сканер nikto для тестирования
веб-приложений.


\printbibliography



\end{document}
